\section{Literature Review}
\label{sec:literature}

\subsection{Dial-a-Ride and Demand-Responsive Transit}

The Dial-a-Ride Problem (DARP) is a variant of the vehicle routing problem with time windows in which a fleet of vehicles must transport passengers between individual origin--destination pairs while respecting time window constraints, maximum ride time limits, and vehicle capacity \citep{cordeau2007}. Since its formal introduction, the DARP has attracted sustained attention as a model for paratransit, elderly and disabled transport, and, more recently, on-demand mobility services in low-density urban areas \citep{psaraftis1988, parragh2008a}.

Demand-responsive transit (DRT) extends the DARP framework to a broader operational context in which service is triggered by real-time passenger requests rather than fixed schedules. The growth of smartphone-based booking platforms and GPS-enabled fleets has made DRT increasingly viable as a first- and last-mile solution in suburban and peri-urban areas where fixed-route bus services are economically unsustainable \citep{molenbruch2017}. Comprehensive surveys by \citet{ho2018} and \citet{pillac2013} document the rapid expansion of DRT research, covering exact methods, metaheuristics, and online algorithms for both static and dynamic problem variants.

A defining characteristic of classical DARP formulations is the assumption of door-to-door service: vehicles travel directly to each passenger's origin and destination address. While this maximizes passenger convenience, it generates highly fragmented routes with significant detour and deadhead mileage, particularly when demand is spatially dispersed \citep{agatz2012}. The present paper relaxes the door-to-door assumption by introducing bidirectional meeting point sets, allowing both pickup and dropoff locations to be selected from a pre-defined network of accessible stops. This structural change is the central departure from classical DARP and motivates the literature reviewed in the following subsection.

\subsection{Meeting Points and Virtual Stops}

The concept of meeting points — fixed or semi-fixed locations where passengers walk a short distance to board or alight — was introduced into the ridesharing literature by \citet{stiglic2015}, who demonstrated that allowing passengers to walk to nearby virtual stops substantially reduces vehicle detour and increases the number of feasible matches in a shared-ride system. A follow-up study \citep{stiglic2018} quantified the trade-off between walking burden and system efficiency, showing that even modest walking tolerances (200--500 metres) can yield significant reductions in total vehicle kilometres travelled without materially degrading passenger experience.

Building on this foundation, \citet{cortenbach2024} introduced the Dial-a-Ride Problem with Meeting Points (DARPmp), the first formal DARP extension to incorporate meeting point selection as a decision variable. Their formulation allows the pickup location to be chosen from a candidate set while the dropoff location remains fixed at the passenger's destination address — a \emph{single-sided} meeting point design. The problem is solved with a Tabu Search metaheuristic applied to a static batch of requests. \citet{cortenbach2024} report substantial efficiency gains over door-to-door baselines, confirming the operational value of flexible pickup locations in many-to-many DRT.

However, the DARPmp of \citet{cortenbach2024} has two important limitations that the present paper addresses. First, the single-sided design leaves dropoff flexibility unexploited: passengers must still be delivered to their exact destination, which constrains route consolidation on the egress leg. Second, the model treats passenger acceptance as deterministic — once the system assigns a meeting point, the passenger is assumed to comply regardless of walking distance or service quality. In practice, passengers exercise choice, and a model that ignores this risks over-optimistic efficiency estimates. The present paper extends DARPmp to a \emph{bidirectional} design in which both the pickup meeting point $m_r^P \in M_r^P$ and the dropoff meeting point $m_r^D \in M_r^D$ are decision variables, and couples this with an explicit multinomial logit (MNL) passenger choice model.

\citet{fielbaum2021} studied bidirectional walking flexibility in ridepooling, showing
that allowing passengers to walk to both pickup and dropoff points reduces vehicle detours
and improves system efficiency. Their analysis demonstrates that the efficiency gains from
dropoff-side flexibility are comparable in magnitude to those from pickup-side flexibility,
providing empirical motivation for the bidirectional design adopted in the present paper.
Our work extends this insight to the many-to-many DRT setting with explicit binary logit passenger
choice and online re-optimization: unlike \citet{fielbaum2021}, who study a static
ridepooling context, we address the online assignment problem where meeting points must be
selected in real time as requests arrive, and passenger acceptance is modeled
probabilistically rather than assumed.

\subsection{Passenger Choice in Demand-Responsive Transit}

Discrete choice modelling provides the theoretical foundation for representing heterogeneous passenger preferences in transportation systems. The multinomial logit (MNL) model, derived from random utility maximisation \citep{benakiva1985}, is the workhorse of mode choice analysis and has been applied extensively to transit demand, ridesharing adoption, and on-demand mobility \citep{lavieri2019}. In the MNL framework, each passenger selects the alternative that maximizes their perceived utility, with choice probabilities determined by the softmax function over deterministic utility components.

The application of discrete choice models to DRT optimization is relatively recent. Most DRT scheduling papers treat passenger acceptance as a binary constraint: a request is either served within its time window or rejected, with no intermediate probability of acceptance. This deterministic treatment ignores the empirical reality that passengers differ in their sensitivity to walking distance, waiting time, in-vehicle travel time, and fare \citep{lavieri2019}. Heterogeneous preferences are particularly consequential in meeting-point DRT, where the system must select not only \emph{whether} to serve a request but \emph{where} to pick up and drop off the passenger — a choice that directly affects the utility experienced.

The present paper builds on the binary logit framework developed in \citet{jin2024work1} for a many-to-one DRT setting with dynamic pricing, extended here to the many-to-many bidirectional meeting-point setting. The core model retains three heterogeneous passenger types and the outside option (passengers may reject all offered bundles). The outside option is critical: it ensures that the optimization does not assign meeting points so inconvenient that passengers would rationally decline service, a failure mode absent from deterministic DARP formulations.

\subsection{Dynamic Scheduling and Rolling Horizon Methods}

Real-world DRT operates in an online environment: requests arrive continuously, vehicles are in motion, and schedules must be updated in near-real-time. The online DARP \citep{psaraftis1988} and its dynamic variants have been studied extensively, with solution approaches ranging from greedy insertion heuristics to population-based metaheuristics and reinforcement learning \citep{pillac2013, ho2018}.

Rolling horizon (receding horizon) methods decompose the infinite-horizon online problem into a sequence of finite-horizon subproblems, each solved to optimality or near-optimality before the horizon advances \citep{bent2004}. \citet{bent2004} established the theoretical basis for rolling horizon approaches in vehicle routing, showing that re-optimization over a sliding window can closely approximate clairvoyant offline solutions when the window length is calibrated to the request arrival rate. This insight motivates the rolling horizon architecture adopted in the present paper.

The Adaptive Large Neighborhood Search (ALNS) of \citet{ropke2006} is the dominant metaheuristic for DARP and related problems, combining destroy-and-repair operators with an adaptive weight mechanism that favors operators performing well on recent iterations. ALNS has been applied to static DARP \citep{parragh2008a}, dynamic ridesharing \citep{agatz2012}, and, most recently, to online DRT re-optimization.

\citet{wu2025} present a rolling horizon ALNS framework for many-to-many DRT that re-optimises vehicle routes at fixed time intervals as new requests arrive. Their computational experiments on synthetic and real-world instances demonstrate that rolling horizon ALNS achieves near-optimal acceptance rates with tractable computation times. However, the formulation of \citet{wu2025} assumes door-to-door service — no meeting points are considered — and treats passenger acceptance as deterministic. The present paper integrates the rolling horizon ALNS architecture of \citet{wu2025} with the bidirectional meeting point assignment and binary logit choice model described above, yielding a unified online co-optimization framework.

\subsection{Research Gap and Positioning}

Table~\ref{tab:literature-gap} summarizes the positioning of this paper relative to the most closely related prior works. To our knowledge, this is the first work to jointly address bidirectional meeting point sets ($M_r^P$ and $M_r^D$), stochastic passenger choice via binary logit with an outside option, rolling horizon re-optimization, and a many-to-many DRT service structure. \citet{fielbaum2021} studied bidirectional walking in static ridepooling; our contribution extends this to online many-to-many DRT with passenger choice and rolling horizon re-optimization.

\begin{table}[h]
\centering
\caption{Positioning relative to prior work}
\label{tab:literature-gap}
\begin{tabular}{lcccc}
\toprule
Study & Bidirectional MPs & Passenger Choice & Rolling Horizon & Many-to-Many \\
\midrule
Cortenbach et al.\ (2024) & No (single-sided) & No & No & Yes \\
Wu et al.\ (2025)         & No                & No & Yes & Yes \\
Fielbaum et al.\ (2021)   & Yes (ridepooling) & No & No & Partial \\
Alonso-Mora et al.\ (2017)& No                & No & No & Yes \\
This paper                & Yes               & Yes (binary logit) & Yes & Yes \\
\bottomrule
\end{tabular}
\end{table}
